
\SGAStatement{Corollary 3.4}{Let \(k\) be a field, and let \(G\) be a
unipotent \(k\)-algebraic group acting on a finite-dimensional
\(k\)-vector space \(V\).  Then \(V\) has a sequence of vector
subspaces \(V_i\), defined over \(k\) and stable under \(G\),
\[
  0=V_0\subset V_1\subset\cdots\subset V_n=V,
\]
such that \(G\) acts trivially on \(V_{i+1}/V_i\).  Moreover, the
\(V_i\) may be chosen so that \(V_{i+1}/V_i\) has dimension \(1\).}
\SGARefTarget{sga3:xvii:corollary:3:4}

We shall now summarize and supplement the properties of unipotent
groups already proved in the following theorem.

\SGAStatement{\SGAInlineRefTarget{sga3:xvii:theorem:3-5}Theorem 3.5}{Let \(G\) be an algebraic group defined
over a field \(k\).  The following properties are equivalent:
\begin{enumerate}
\item[(i)] \SGARefTarget{sga3:xvii:theorem:3:5:item:i}
  \(G\) is unipotent.
\item[(ii)] \SGARefTarget{sga3:xvii:theorem:3:5:item:ii}
  \(G\) has a composition series defined over \(k\), whose successive
  quotients are isomorphic to \(\mathbb G_a\) if \(p=0\), and,
  respectively, to \(\boldsymbol\alpha_p\), \(\mathbb G_a\), or a
  twisted form of \((\mathbb Z/p\mathbb Z)^r\) (\SGARefLink{sga3:xvii:definition:1:4}{1.4}) if \(p>0\).
\item[(iii)] \SGARefTarget{sga3:xvii:theorem:3:5:item:iii}
  As in (ii), with the additional requirement that the composition
  series be central.
\item[(iv)] \SGARefTarget{sga3:xvii:theorem:3:5:item:iv}
  \(G\) has a characteristic composition series
  (Exposé VI\textsubscript{B}) defined over \(k\), whose successive
  quotients are isomorphic to \((\mathbb G_a)^r\) if \(p=0\), and,
  respectively, to \((\boldsymbol\alpha_p)^r\), a twisted form of
  \((\mathbb G_a)^s\), or a twisted form of
  \((\mathbb Z/p\mathbb Z)^t\), taken in this order, if \(p>0\).
\item[(v)] \SGARefTarget{sga3:xvii:theorem:3:5:item:v}
  \(G\) is isomorphic to an algebraic subgroup of the group
  \(\operatorname{Trigstr}(n)_k\) of strictly upper-triangular
  matrices in the linear group \(\operatorname{GL}(n)_k\), for a
  suitable integer \(n\geq0\).
\item[(vi)] \SGARefTarget{sga3:xvii:theorem:3:5:item:vi}
  \(G\) is affine, and for every linear representation of \(G\) on a
  nonzero finite-dimensional \(k\)-vector space \(V\), one has
  \(V^G\ne0\).
\end{enumerate}}
\SGARefTarget{sga3:xvii:theorem:3:5}

\SGAProof{Proof}
The implication \SGARefLink{sga3:xvii:theorem:3:5:item:i}{(i)}\(\Rightarrow\)\SGARefLink{sga3:xvii:theorem:3:5:item:vi}{(vi)} follows from \SGARefLink{sga3:xvii:proposition:2:1}{2.1} and \SGARefLink{sga3:xvii:proposition:3:2}{3.2}.

For \SGARefLink{sga3:xvii:theorem:3:5:item:vi}{(vi)}\(\Rightarrow\)\SGARefLink{sga3:xvii:theorem:3:5:item:v}{(v)}, since the algebraic group \(G\) is
affine, it is an algebraic subgroup of a suitable linear group
\(\operatorname{GL}(V)\) (Exposé VI\textsubscript{B}, \SGARefLink{def:sga3-VIB-11.3}{11.3}).
Applying \SGARefLink{sga3:xvii:corollary:3:4}{3.4} to the representation of \(G\) on \(V\) defined by this
embedding gives \SGARefLink{sga3:xvii:theorem:3:5:item:v}{(v)}.

For \SGARefLink{sga3:xvii:theorem:3:5:item:v}{(v)}\(\Rightarrow\)\SGARefLink{sga3:xvii:theorem:3:5:item:iii}{(iii)}, the algebraic group
\(\operatorname{Trigstr}(n)\) is known to have a central composition
series whose successive quotients are isomorphic to \(\mathbb G_a\).
The composition series induced on \(G\) gives property \SGARefLink{sga3:xvii:theorem:3:5:item:iii}{(iii)}, in view
of \SGARefLink{sga3:xvii:proposition:1:5}{1.5}.

The implications
\[
  \text{(iii)}\Rightarrow\text{(ii)}\Rightarrow\text{(i)}
  \qquad\text{and}\qquad
  \text{(iv)}\Rightarrow\text{(i)}
\]
are clear.

We shall prove \SGARefLink{sga3:xvii:theorem:3:5:item:i}{(i)}\(\Rightarrow\)\SGARefLink{sga3:xvii:theorem:3:5:item:iv}{(iv)} shortly, but first note some
consequences of what has already been proved.

\SGAStatement{Definition 3.6}{A \(p\)-Lie algebra \(\mathfrak g\),
\(p>0\) (cf. Exposé VII\textsubscript{A}, \S5), is called
\emph{unipotent} if the map
\[
  x\longmapsto x^{(p)}
\]
is nilpotent; that is, if for every \(x\in\mathfrak g\) there is an
integer \(n>0\) such that \(x^{(p^n)}=0\).}
\SGARefTarget{sga3:xvii:definition:3:6}
